\hypertarget{huawei-ai-certification-lab-9.5}{%
\section{Huawei AI Certification: Lab
9.5}\label{huawei-ai-certification-lab-9.5}}

\hypertarget{modelarts-exeml-equivalent-local-image-classification}{%
\subsection{ModelArts ExeML Equivalent (Local Image
Classification)}\label{modelarts-exeml-equivalent-local-image-classification}}

\textbf{Objective:} Construct an Image Classification pipeline
replicating the ModelArts capability locally. By leveraging Transfer
Learning (MobileNetV2), we systematically ingest massive real-world
Image Datasets (flowers) and abstract advanced feature arrays accurately
and highly computationally efficient!

\begin{lstlisting}[language=Python]
import pathlib
import tensorflow as tf
from tensorflow.keras import layers, models
import matplotlib.pyplot as plt
import numpy as np

%matplotlib inline
\end{lstlisting}

\begin{lstlisting}
WARNING: All log messages before absl::InitializeLog() is called are written to STDERR
I0000 00:00:1777684688.089842 2694181 cudart_stub.cc:31] Could not find cuda drivers on your machine, GPU will not be used.
I0000 00:00:1777684688.919096 2694181 cpu_feature_guard.cc:227] This TensorFlow binary is optimized to use available CPU instructions in performance-critical operations.
To enable the following instructions: AVX2 FMA, in other operations, rebuild TensorFlow with the appropriate compiler flags.
WARNING: All log messages before absl::InitializeLog() is called are written to STDERR
I0000 00:00:1777684690.385648 2694181 cudart_stub.cc:31] Could not find cuda drivers on your machine, GPU will not be used.
\end{lstlisting}

\hypertarget{cloud-dataset-acquisition-extraction}{%
\subsubsection{1. Cloud Dataset Acquisition \&
Extraction}\label{cloud-dataset-acquisition-extraction}}

Programmatically initiating internal retrieval blocks across HTTPS to
safely extract raw Image folders into classification directories.

\begin{lstlisting}[language=Python]
PROJECT_DIR = pathlib.Path.cwd()
DATASET_DIR = PROJECT_DIR / 'data'
DATASET_DIR.mkdir(exist_ok=True)

dataset_url = "https://storage.googleapis.com/download.tensorflow.org/example_images/flower_photos.tgz"
archive = tf.keras.utils.get_file(
    'flower_photos.tgz',
    origin=dataset_url,
    extract=True,
    cache_dir=str(DATASET_DIR),
    cache_subdir='',
)
archive_path = pathlib.Path(archive)
candidates = [
    archive_path.parent / 'flower_photos_extracted' / 'flower_photos',
    archive_path.parent / 'flower_photos',
]
data_dir = next((p for p in candidates if p.exists()), None)
if data_dir is None:
    raise FileNotFoundError(f"Missing flower_photos beside {archive_path}")
print("Dataset securely mounted at:", data_dir)
\end{lstlisting}

\begin{lstlisting}
Downloading data from https://storage.googleapis.com/download.tensorflow.org/example_images/flower_photos.tgz
228813984/228813984 ━━━━━━━━━━━━━━━━━━━━ 7s 0us/step
Dataset securely mounted at: /home/ansinitro/AITU/ai-nn/midterm/Labs_notebooks/9.5 ModelArts Lab Guide/data/flower_photos_extracted/flower_photos
\end{lstlisting}

\hypertarget{generative-image-loaders-validations}{%
\subsubsection{2. Generative Image Loaders \&
Validations}\label{generative-image-loaders-validations}}

Using native \passthrough{\lstinline!image\_dataset\_from\_directory!}
to mathematically divide training partitions correctly and batching the
vectors.

\begin{lstlisting}[language=Python]
IMG_SIZE = 180
BATCH_SIZE = 32
EPOCHS = 10

train_ds = tf.keras.utils.image_dataset_from_directory(
    data_dir,
    validation_split=0.2,
    subset="training",
    seed=123,
    image_size=(IMG_SIZE, IMG_SIZE),
    batch_size=BATCH_SIZE
)

val_ds = tf.keras.utils.image_dataset_from_directory(
    data_dir,
    validation_split=0.2,
    subset="validation",
    seed=123,
    image_size=(IMG_SIZE, IMG_SIZE),
    batch_size=BATCH_SIZE
)

class_names = train_ds.class_names
print("Target Classification Classes:", class_names)
\end{lstlisting}

\begin{lstlisting}
Found 3670 files belonging to 5 classes.
Using 2936 files for training.
Found 3670 files belonging to 5 classes.
Using 734 files for validation.
Target Classification Classes: ['daisy', 'dandelion', 'roses', 'sunflowers', 'tulips']


E0000 00:00:1777684701.989102 2694181 cuda_executor.cc:1737] INTERNAL: CUDA Runtime error: Failed call to cudaGetRuntimeVersion: Error loading CUDA libraries. GPU will not be used.: Error loading CUDA libraries. GPU will not be used.
W0000 00:00:1777684701.989408 2694836 cuda_executor.cc:1755] Failed to determine cuDNN version (Note that this is expected if the application doesn't link the cuDNN plugin): INTERNAL: cuDNN error: CUDNN_STATUS_INTERNAL_ERROR
W0000 00:00:1777684702.016392 2694181 gpu_device.cc:2365] Cannot dlopen some GPU libraries. Please make sure the missing libraries mentioned above are installed properly if you would like to use GPU. Follow the guide at https://www.tensorflow.org/install/gpu for how to download and setup the required libraries for your platform.
Skipping registering GPU devices...
\end{lstlisting}

\begin{lstlisting}[language=Python]
plt.figure(figsize=(10, 10))
for images, labels in train_ds.take(1):
    for i in range(9):
        plt.subplot(3, 3, i + 1)
        plt.imshow(images[i].numpy().astype("uint8"))
        plt.title(class_names[labels[i]])
        plt.axis("off")
plt.suptitle('Dataset Class Sample Visualization')
plt.tight_layout()
plt.show()
\end{lstlisting}

\begin{figure}[H]
\centering
\includegraphics[width=0.85\linewidth]{"9.5 ModelArts Lab Guide/9.5 ModelArts Lab Guide_files/9.5 ModelArts Lab Guide_6_0.png"}
\end{figure}

\hypertarget{transfer-learning-mobilenetv2-architecture-pipeline}{%
\subsubsection{3. Transfer Learning (MobileNetV2) Architecture
Pipeline}\label{transfer-learning-mobilenetv2-architecture-pipeline}}

Freezing the computationally heavy \passthrough{\lstinline!imagenet!}
parameters, mapping our arbitrary classes via pooling structures
seamlessly.

\begin{lstlisting}[language=Python]
AUTOTUNE = tf.data.AUTOTUNE
train_ds = train_ds.cache().shuffle(1000).prefetch(buffer_size=AUTOTUNE)
val_ds = val_ds.cache().prefetch(buffer_size=AUTOTUNE)

base_model = tf.keras.applications.MobileNetV2(
    input_shape=(IMG_SIZE, IMG_SIZE, 3),
    include_top=False,
    weights='imagenet'
)
base_model.trainable = False

model = models.Sequential([
    layers.Rescaling(1./255, input_shape=(IMG_SIZE, IMG_SIZE, 3)),
    base_model,
    layers.GlobalAveragePooling2D(),
    layers.Dense(128, activation='relu'),
    layers.Dropout(0.2),
    layers.Dense(len(class_names), activation='softmax')
])

model.compile(
    optimizer='adam',
    loss='sparse_categorical_crossentropy',
    metrics=['accuracy']
)

print("Architecture successfully locked! Proceeding to execution pipeline...")
\end{lstlisting}

\begin{lstlisting}
/tmp/ipykernel_2694181/1128898239.py:5: UserWarning: `input_shape` is undefined or non-square, or `rows` is not in [96, 128, 160, 192, 224]. Weights for input shape (224, 224) will be loaded as the default.
  base_model = tf.keras.applications.MobileNetV2(


Downloading data from https://storage.googleapis.com/tensorflow/keras-applications/mobilenet_v2/mobilenet_v2_weights_tf_dim_ordering_tf_kernels_1.0_224_no_top.h5
9406464/9406464 ━━━━━━━━━━━━━━━━━━━━ 1s 0us/step
Architecture successfully locked! Proceeding to execution pipeline...


/home/ansinitro/AITU/ai-nn/midterm/.venv/lib/python3.12/site-packages/keras/src/layers/preprocessing/data_layer.py:95: UserWarning: Do not pass an `input_shape`/`input_dim` argument to a layer. When using Sequential models, prefer using an `Input(shape)` object as the first layer in the model instead.
  super().__init__(**kwargs)
\end{lstlisting}

\hypertarget{evaluation-loop-graphical-validation}{%
\subsubsection{4. Evaluation Loop \& Graphical
Validation}\label{evaluation-loop-graphical-validation}}

Deploy the execution across configured Epochs and visualize structural
accuracy improvements natively.

\begin{lstlisting}[language=Python]
history = model.fit(train_ds, validation_data=val_ds, epochs=EPOCHS)

fig, axes = plt.subplots(1, 2, figsize=(12, 5))
axes[0].plot(history.history['accuracy'], label='Train')
axes[0].plot(history.history['val_accuracy'], label='Val')
axes[0].set_xlabel('Epoch')
axes[0].set_ylabel('Accuracy')
axes[0].set_title('Accuracy Trajectory')
axes[0].legend()
axes[0].grid(alpha=0.3)

axes[1].plot(history.history['loss'], label='Train')
axes[1].plot(history.history['val_loss'], label='Val')
axes[1].set_xlabel('Epoch')
axes[1].set_ylabel('Loss')
axes[1].set_title('Loss Optimization Trajectory')
axes[1].legend()
axes[1].grid(alpha=0.3)

plt.suptitle('Flower Classification Optimization')
plt.tight_layout()
plt.show()
\end{lstlisting}

\begin{lstlisting}
Epoch 1/10


W0000 00:00:1777684711.455531 2694844 cpu_allocator_impl.cc:82] Allocation of 99532800 exceeds 10% of free system memory.
W0000 00:00:1777684711.610671 2694844 cpu_allocator_impl.cc:82] Allocation of 101756928 exceeds 10% of free system memory.
W0000 00:00:1777684711.865218 2694844 cpu_allocator_impl.cc:82] Allocation of 37324800 exceeds 10% of free system memory.
W0000 00:00:1777684711.888296 2694844 cpu_allocator_impl.cc:82] Allocation of 37324800 exceeds 10% of free system memory.
W0000 00:00:1777684711.925538 2694844 cpu_allocator_impl.cc:82] Allocation of 37324800 exceeds 10% of free system memory.


92/92 ━━━━━━━━━━━━━━━━━━━━ 56s 542ms/step - accuracy: 0.7807 - loss: 0.6213 - val_accuracy: 0.8692 - val_loss: 0.4010
Epoch 2/10
92/92 ━━━━━━━━━━━━━━━━━━━━ 46s 502ms/step - accuracy: 0.8811 - loss: 0.3250 - val_accuracy: 0.8719 - val_loss: 0.3708
Epoch 3/10
92/92 ━━━━━━━━━━━━━━━━━━━━ 81s 497ms/step - accuracy: 0.9135 - loss: 0.2391 - val_accuracy: 0.8828 - val_loss: 0.3387
Epoch 4/10
92/92 ━━━━━━━━━━━━━━━━━━━━ 45s 490ms/step - accuracy: 0.9387 - loss: 0.1812 - val_accuracy: 0.8856 - val_loss: 0.3675
Epoch 5/10
92/92 ━━━━━━━━━━━━━━━━━━━━ 47s 511ms/step - accuracy: 0.9564 - loss: 0.1228 - val_accuracy: 0.8760 - val_loss: 0.3412
Epoch 6/10
92/92 ━━━━━━━━━━━━━━━━━━━━ 52s 570ms/step - accuracy: 0.9762 - loss: 0.0861 - val_accuracy: 0.8801 - val_loss: 0.3552
Epoch 7/10
92/92 ━━━━━━━━━━━━━━━━━━━━ 76s 509ms/step - accuracy: 0.9830 - loss: 0.0660 - val_accuracy: 0.8856 - val_loss: 0.3539
Epoch 8/10
92/92 ━━━━━━━━━━━━━━━━━━━━ 45s 494ms/step - accuracy: 0.9877 - loss: 0.0521 - val_accuracy: 0.8787 - val_loss: 0.3625
Epoch 9/10
92/92 ━━━━━━━━━━━━━━━━━━━━ 45s 495ms/step - accuracy: 0.9911 - loss: 0.0410 - val_accuracy: 0.8883 - val_loss: 0.3527
Epoch 10/10
92/92 ━━━━━━━━━━━━━━━━━━━━ 46s 497ms/step - accuracy: 0.9918 - loss: 0.0376 - val_accuracy: 0.8774 - val_loss: 0.3695
\end{lstlisting}

\begin{figure}[H]
\centering
\includegraphics[width=0.85\linewidth]{"9.5 ModelArts Lab Guide/9.5 ModelArts Lab Guide_files/9.5 ModelArts Lab Guide_10_3.png"}
\end{figure}

\hypertarget{prediction-confusion-matrix}{%
\subsubsection{5. Prediction Confusion
Matrix}\label{prediction-confusion-matrix}}

Isolating classification false-positive intersections structurally.

\begin{lstlisting}[language=Python]
all_true = []
all_pred = []
for imgs, lbls in val_ds:
    p = model.predict(imgs, verbose=0)
    all_pred.append(np.argmax(p, axis=1))
    all_true.append(lbls.numpy())
all_true = np.concatenate(all_true)
all_pred = np.concatenate(all_pred)

cm = tf.math.confusion_matrix(all_true, all_pred, num_classes=len(class_names)).numpy()

fig, ax = plt.subplots(figsize=(7, 6))
im = ax.imshow(cm, cmap='Blues')
ax.set_xticks(range(len(class_names)))
ax.set_yticks(range(len(class_names)))
ax.set_xticklabels(class_names, rotation=45, ha='right')
ax.set_yticklabels(class_names)
ax.set_xlabel('Predicted Classes')
ax.set_ylabel('True Classes')
ax.set_title('Model validation spatial intersection map')
threshold = cm.max() / 2
for i in range(cm.shape[0]):
    for j in range(cm.shape[1]):
        ax.text(j, i, str(cm[i, j]), ha='center', va='center', color='white' if cm[i, j] > threshold else 'black')
plt.colorbar(im, ax=ax)
plt.tight_layout()
plt.show()
\end{lstlisting}

\begin{figure}[H]
\centering
\includegraphics[width=0.85\linewidth]{"9.5 ModelArts Lab Guide/9.5 ModelArts Lab Guide_files/9.5 ModelArts Lab Guide_12_0.png"}
\end{figure}

\hypertarget{endpoint-prediction-execution}{%
\subsubsection{6. Endpoint Prediction
Execution}\label{endpoint-prediction-execution}}

\begin{lstlisting}[language=Python]
images_batch, labels_batch = next(iter(val_ds))
preds = model.predict(images_batch, verbose=0)
pred_labels = np.argmax(preds, axis=1)
true_labels = labels_batch.numpy()

n_show = min(9, len(images_batch))
correct_total = int((pred_labels == true_labels).sum())

plt.figure(figsize=(10, 10))
for i in range(n_show):
    plt.subplot(3, 3, i + 1)
    plt.imshow(images_batch[i].numpy().astype("uint8"))
    is_correct = pred_labels[i] == true_labels[i]
    color = 'green' if is_correct else 'red'
    mark = '✓' if is_correct else '✗'
    plt.title(f"{mark} pred: {class_names[pred_labels[i]]}\nreal: {class_names[true_labels[i]]}", color=color, fontsize=10)
    plt.axis("off")
plt.suptitle(f'Verification Results on Validation Batch: {correct_total}/{len(images_batch)} Accurate')
plt.tight_layout()
plt.show()
\end{lstlisting}

\begin{figure}[H]
\centering
\includegraphics[width=0.85\linewidth]{"9.5 ModelArts Lab Guide/9.5 ModelArts Lab Guide_files/9.5 ModelArts Lab Guide_14_0.png"}
\end{figure}
